\chapter{Literature Review}

\section{Introduction to Concurrency and Transactional Memory}

In concurrent systems, shared state must be protected to ensure correctness.

Dijkstra~\cite{dijkstra1968} introduced semaphores as a low-level primitive for protecting
critical sections. A semaphore is a non-negative integer controlled by two atomic operations:
P (\emph{passeren}, to pass) and V (\emph{vrijgeven}, to release). The P operation decrements
the semaphore when its value is positive and otherwise blocks, while the V operation increments
the semaphore and may wake a waiting thread. Because P and V are independent operations, with
nothing enforcing that each acquisition is matched by a respective release, it is easy to
forget one.

In addition, an exit from a critical section may require releasing all acquired
semaphores; failing to do so causes deadlock, as illustrated in Listing~\ref{lst:semaphore}.

\begin{listing}[ht]
  \begin{minted}{clojure}
(defn transfer [from to amount]
  (.acquire (:lock from))
  (.acquire (:lock to))
  ;; Exception occurring after lock acquisition
  ;; prevents the locks from being released.
  (when (<= amount 0)
    (throw (Exception. "Invalid transfer amount")))
  (when (< (:balance from) amount)
    (throw (Exception. "Insufficient funds")))
  (set! (:balance from) 
        (- (:balance from) amount))
  (set! (:balance to)   
        (+ (:balance to) amount))
  (.release (:lock to))
  (.release (:lock from)))
\end{minted}
  \caption{Explicit lock-based synchronization vulnerable to lock leakage under exceptions.}
  \label{lst:semaphore}
\end{listing}

To address these limitations, Hoare introduced monitors~\cite{hoare1974}.
A monitor operates on a single lock that is automatically acquired upon entry and released
upon every exit, whether normal or exceptional.

\begin{listing}[ht]
  \begin{minted}{clojure}
(defn transfer-safe [from to amount]
  (locking from
    (locking to
      (when (<= amount 0)
        (throw (Exception. "Invalid transfer amount")))
      (when (< (:balance from) amount)
        (throw (Exception. "Insufficient funds")))
      (set! (:balance from) 
            (- (:balance from) amount))
      (set! (:balance to)   
            (+ (:balance to) amount)))))
\end{minted}
  \caption{Monitor-style synchronization with automatic lock release.}
  \label{lst:monitor}
\end{listing}

However, monitors do not eliminate concurrency challenges. Coarse-grained
locking\footnote{A synchronization strategy where a single lock
  protects a large portion of data.} reduces
parallelism and provides no rollback semantics. If an exception occurs after
partial state modification, those intermediate updates remain visible unless
manually undone by the programmer.

These limitations motivated the search for mechanisms that allow programmers to express atomicity declaratively.
Lomet~\cite{lomet1977} proposed that programmers should be
able to specify that a sequence of operations must execute atomically without explicitly
managing synchronization details, an idea that was later realized in transactional
memory systems.

\subsection{Hardware Transactional Memory}

Herlihy and Moss~\cite{herlihy1993} introduced
Transactional Memory (TM) as a hardware-level mechanism for atomic operations
without explicit locking. A transaction is a finite sequence of machine
instructions satisfying serializability and atomicity~\cite{herlihy1993}.
HTM exposes five primitive instructions:

\begin{itemize}
  \item \textsc{Load-transactional} (LT) reads a shared memory location into
        a private register.
  \item \textsc{Load-transactional-exclusive} (LTX) reads a shared memory
        location into a private register, hinting it will be updated.
  \item \textsc{Store-transactional} (ST) speculatively writes to a shared
        location; the write remains invisible until the transaction commits.
  \item \textsc{Validate} (VALIDATE) checks whether the read set remains
        consistent; returns \emph{False} and discards tentative updates if
        another transaction has intervened.
  \item \textsc{Commit} (COMMIT) attempts to make all tentative writes
        permanent, succeeding only if no conflicting transaction has modified
        the data set.
\end{itemize}

Rather than acquiring a lock, a process would: (1) use LT or LTX to read
locations, (2) \textsc{Validate} consistency, (3) use ST to speculatively
write, and (4) \textsc{Commit}. On failure at any step the process retries,
possibly applying \emph{adaptive backoff} under high
contention~\cite{herlihy1993}.

\begin{listing}[ht]
  \begin{minted}{clojure}
(defn transfer-htm [account-a account-b amount]
  (loop []
    (let [bal-a (LT account-a)
          bal-b (LT account-b)]
      (if (VALIDATE)
        (do
          (ST account-a (- bal-a amount))
          (ST account-b (+ bal-b amount))
          (if (COMMIT)
            :success
            (do (adaptive-backoff!) (recur))))
        (do (adaptive-backoff!) (recur))))))
\end{minted}
  \caption{HTM-style transaction.}
\end{listing}

Two conflict classes can cause \textsc{Validate} or \textsc{Commit} to fail:

\begin{table}[h]
  \centering
  \begin{tabular}{ll}
    \hline
    \textbf{Conflict Type} & \textbf{Example}                       \\
    \hline
    Read-Write             & $T_1$ reads $A$;\ $T_2$ commits $W(A)$ \\
    Write-Write            & $T_1$ and $T_2$ both attempt $W(B)$    \\
    \hline
  \end{tabular}
  \caption{Transaction conflict classes illustrated by the transfer example.}
\end{table}

\subsection{Software Transactional Memory}

Shavit and Touitou~\cite{shavit1995} first introduced a software transactional
memory model and a corresponding implementation algorithm without hardware support,
establishing the conceptual foundation for subsequent STM research. However, their
approach was restricted to \emph{static transactions}.%
\footnote{Transactions which access a pre-determined sequence of locations.}
\citeauthor{dice2006transactional} observe that the dominant trend in
transactional memory design has been toward transactions that are
\emph{unbounded}, placing no limit on the number of locations a transaction
may access, and \emph{dynamic}, meaning those locations are determined only
at runtime rather than fixed in advance~\cite{dice2006transactional}.

\begin{listing}[ht]
  \begin{minted}{clojure}
; Static: data declared before execution
(defn transfer [account-a account-b amount]
  (dosync
    (alter account-a - amount)
    (alter account-b + amount)))

; Dynamic: data discovered at runtime
(dosync
  (let [node @head]
    (ref-set head @(:next node))))
\end{minted}
  \caption{Static vs.\ dynamic transactions.}
\end{listing}

Since the data must be known statically, any transaction that discovers its
memory locations at runtime (such as traversing a linked list or walking a tree)
cannot be expressed in Shavit and Touitou's model. This restriction motivated later
STM designs. It is also directly relevant to our system: Clojure programs regularly
discover their data at runtime, selecting accounts at random as in the bank
transfer benchmark, or following a chain of \texttt{deref} calls through dynamic
data structures. Static transactions are therefore an unsuitable foundation,
and the system we present is fully dynamic.

\subsection{Concurrency Control}

Transactional Memory systems must ensure that concurrent transactions produce
results equivalent to a serial execution. Achieving such serializability requires
resolving conflicts between transactions. Two broad strategies are commonly employed:
pessimistic control, which prevents conflicts through eager lock acquisition,
and optimistic control, which permits speculative execution and detects conflicts
during validation or commit. The following subsections examine these approaches and
their implications for STM design.

\subsubsection{Pessimistic Concurrency Control}

Pessimistic concurrency control assumes that conflicts between transactions are
likely and therefore prevents them by acquiring locks before shared data is
accessed~\cite{bernstein1987}. Consider two concurrent transactions $T_1$ and
$T_2$ updating shared variables $A$ and $B$:

\begin{itemize}
  \item $T_1$: $R(A)\ W(B)$
  \item $T_2$: $W(A)\ W(B)$
\end{itemize}

These operations introduce read-write ($R(A) \times W(A)$) and write-write
($W(B) \times W(B)$) conflicts. Strict two-phase locking (S2PL) prevents such
conflicts by requiring shared locks (s-locks) for reads and exclusive locks
(x-locks) for writes, which are retained until commit or rollback~\cite{bernstein1987}.
If $T_1$ acquires an s-lock on $A$, $T_2$ cannot obtain the x-lock required to
modify $A$ until $T_1$ completes, and similarly for $B$.

This behavior follows from two fundamental properties of S2PL:

\begin{enumerate}
  \item If a transaction holds an s-lock on a data item, no other active transaction may concurrently modify that item.
  \item If a transaction holds an x-lock on a data item, no other active transaction may concurrently access that item.
\end{enumerate}


Although S2PL guarantees serializability, holding locks until transaction completion reduces concurrency.
Releasing locks early would improve performance but violates serializability. For example, if $T_1$
reads and unlocks $A$ before accessing $B$, $T_2$ may interleave by modifying $B$ and committing in between,
producing a non-serializable schedule.

This blocking behavior motivated optimistic concurrency control, where transactions execute speculatively
and validate only at commit time. Multi-Version Concurrency Control (MVCC) is a realization of this
approach~\cite{bernstein1983} and forms the basis of Clojure’s STM~\cite{clojureRefs}.

\subsubsection{Multi-Version Concurrency Control (MVCC)}

In a multiversion concurrency control, each write on a data
produces a new version rather than overwriting it~\cite{bernstein1983}.
Each version is timestamped with the commit time of the transaction that wrote
it.

Consider a shared object $O$ accessed by three transactions $T_1$, $T_2$,
$T_3$. $T_1$ commits at $t=3$ writing $O=20$, and $T_3$ commits at $t=11$
writing $O=30$. A read-only transaction\footnote{A transaction that performs only read
  operations and does not modify shared state.} $T_2$ begins at $t=5$:

\begin{enumerate}
  \item $T_1$ commits, creating version $v_1$ with $\mathit{t}(v_1) = 3$.
  \item $T_3$ commits, creating version $v_2$ with $\mathit{t}(v_2) = 11$.
  \item $T_2$ begins at $t=5$ and reads $O$. It receives $v_1$ (value $=
          20$), the latest version committed before $t=5$, even though $T_3$
        has already committed $v_2$. $T_2$ obtains a consistent read without
        acquiring any lock.
\end{enumerate}

Since read-only transactions have a consistent view of the data, they never
acquire locks and therefore never have to wait~\cite{bernstein1983}.
The key design implication is that MVCC
separates the concern of read consistency from the concern of write
serialization, that makes MVCC well-suited to read-dominated workloads, and it
is precisely the property exploited in
Section~\ref{sec:read-write-heavy}, where 90\% read-only transactions achieve
the highest throughput of any evaluated scenario by avoiding the commit lock
entirely.

\subsubsection{Transactional Locking II (TL2)}

Transactional Locking II (TL2), proposed by
\citeauthor{dice2006transactional}~\cite{dice2006transactional}, is a software
transactional memory algorithm based on commit-time lock acquisition and a
global version clock. The central insight is that a shared version clock,
incremented once by each writing transaction, allows conflicts to be detected
efficiently without holding locks during speculative
execution~\cite{dice2006transactional}.

Each transaction begins by sampling the global version clock into a local
read-version number $rv$. During speculative execution the transaction records
a \emph{read-set} of accessed locations and a \emph{write-set} of buffered
modifications. Each transactional load is post-validated inline: if the
location's version number exceeds $rv$, the transaction aborts
immediately~\cite{dice2006transactional}. No locks are acquired during this
phase.

At commit time, write transactions acquire versioned write-locks for all
entries in their write-set. TL2 supports two lock-mapping policies: per-object
(PO), where each shared object carries its own lock, and per-stripe (PS),
where a fixed lock array is shared across locations via
hashing~\cite{dice2006transactional}. PS reduces memory overhead but
introduces occasional false conflicts between unrelated locations that hash to
the same stripe. After acquiring locks, the transaction validates its read-set,
increments the global version clock, applies buffered writes, and releases
locks. If validation fails the transaction aborts and retries.
Figure~\ref{fig:tl2-write} illustrates the complete write transaction protocol.

\begin{figure}[htbp]
  \centering
  \begin{tikzpicture}[
    node distance = 0.65cm and 1.4cm,
    font = \sffamily\footnotesize,
    process/.style = {rectangle, minimum width=4.3cm, text width=3.9cm,
        minimum height=0.9cm, text centered, draw=black, thick,
        fill=white, align=center},
    abort/.style = {rectangle, minimum width=2.8cm, text width=2.6cm,
        minimum height=0.8cm, text centered, draw=red!70!black, thick,
        fill=red!5, align=center},
    inline/.style = {rectangle, minimum width=3.4cm, text width=3.0cm,
        minimum height=0.8cm, text centered, draw=orange!70!black, thick,
        fill=orange!5, align=center},
    arrow/.style = {thick, -{Latex[scale=1.0]}},
    darrow/.style = {thick, dashed, draw=orange!70!black, -{Latex[scale=1.0]}}
    ]
    \node (s1) [process] {\textbf{Sample version clock} \\
      \textit{\scriptsize store as $rv$}};
    \node (s2) [process, below=of s1] {\textbf{Speculative execution} \\
      \textit{\scriptsize build read-set, write-set}};
    \node (s3) [process, below=of s2] {\textbf{Lock write-set} \\
      \textit{\scriptsize acquire per-object or stripe locks}};
    \node (s4) [process, below=of s3] {\textbf{Validate read-set} \\
      \textit{\scriptsize all versions $\leq rv$?}};
    \node (s5) [process, below=of s4] {\textbf{Increment version clock} \\
      \textit{\scriptsize store as $wv$}};
    \node (s6) [process, below=of s5] {\textbf{Apply writes + release locks} \\
      \textit{\scriptsize stamp each location with $wv$}};

    \node (chk) [inline, right=of s2] {\textbf{Inline check (per read)} \\
      \textit{\scriptsize version $> rv \Rightarrow$ abort}};
    \node (ab1) [abort, right=of s4]  {\textbf{Abort + retry} \\
      \textit{\scriptsize read-set invalid}};

    \draw [arrow]  (s1) -- (s2);
    \draw [arrow]  (s2) -- (s3);
    \draw [arrow]  (s3) -- (s4);
    \draw [arrow]  (s4) -- node[left]{\textit{pass}} (s5);
    \draw [arrow]  (s5) -- (s6);
    \draw [darrow] (s2) -- (chk);
    \draw [arrow]  (s4) -- node[above]{\textit{fail}} (ab1);
  \end{tikzpicture}
  \caption{TL2 write transaction protocol.}
  \label{fig:tl2-write}
\end{figure}

TL2's separation of read-set and write-set, together with its global
version clock, provides the foundation for versioned data. While
TL2 itself keeps only the most recent version of each item, the
addition of version history would produce a design aligned with
MVCC. The design decisions and extensions adopted in this thesis are
discussed in Chapter~\ref{ch:implementation}.

\subsubsection{Synchronization Anomalies in STM}

Parri~\cite{parriTmIntro} presents a taxonomy of synchronization anomalies in
Software Transactional Memory. Because our STM system is based on optimistic MVCC,
dirty reads are excluded by design.
Two anomalies from Parri's taxonomy remain relevant:

\begin{itemize}
  \item \textbf{Non-Repeatable Read}: a transaction reads the same
        location twice and observes different values.
  \item \textbf{Intermediate Lost Update}: concurrent transactions modify
        overlapping locations and one update is silently overwritten.
\end{itemize}

In our system, non-repeatable reads are addressed
through Multi-Version Concurrency Control (MVCC) and read-set validation,
while intermediate lost updates are prevented through ownership-based conflict
resolution. These mechanisms are discussed further in
Sections~\ref{sec:reading} and~\ref{sec:conflict-resolution}.

\section{Existing Clojure STM Approaches and Limitations}

Clojure's STM exposes a concise API built around \texttt{dosync}, with
a function body, in which reads and writes to \texttt{Refs} are
executed atomically~\cite{clojureRefs}. The implementation is within
the \texttt{LockingTransaction} class.
Table~\ref{tab:clojure-stm} summarizes the key internal components of
\texttt{LockingTransaction} and \texttt{Ref}.

\begin{table}[htbp]
  \centering
  \small
  \begin{tabular}{lll}
    \toprule
    \textbf{Component} & \textbf{Field}      & \textbf{Role}                                     \\
    \midrule
    \multirow{5}{*}{\texttt{LockingTransaction}}
                       & \texttt{startPoint} & Globally unique transaction identifier            \\
                       & \texttt{readPoint}  & Snapshot point for MVCC reads                     \\
                       & \texttt{vals}       & In-transaction buffered write values              \\
                       & \texttt{sets}       & Refs modified via \texttt{alter}/\texttt{ref-set} \\
                       & \texttt{status}     & \texttt{RUNNING}, \texttt{COMMITTING},
    \texttt{COMMITTED}, \texttt{RETRY}, \texttt{KILLED}                                          \\
    \midrule
    \multirow{3}{*}{\texttt{Ref}}
                       & \texttt{tvals}      & Doubly linked history of committed values         \\
                       & \texttt{maxHistory} & Upper bound on retained versions                  \\
                       & \texttt{tinfo}      & Reference to owning transaction's \texttt{Info}   \\
    \bottomrule
  \end{tabular}
  \caption{Key internal fields of Clojure's STM implementation~\cite{jensen2016}.}
  \label{tab:clojure-stm}
\end{table}

A read inside a transaction calls \texttt{doGet}, which acquires a read lock
on the \texttt{Ref} and traverses \texttt{tvals} backwards to find the latest
committed value whose \texttt{point} does not exceed
\texttt{readPoint}~\cite{jensen2016}. If no such value exists, a fault counter
is incremented and the transaction retries, triggering an expansion of the
\texttt{Ref}'s history on the next attempt. A write calls \texttt{doSet},
which acquires a write lock on the \texttt{Ref}, adds it to \texttt{sets}, and
buffers the new value in \texttt{vals} without modifying shared memory until
commit~\cite{jensen2016}.

However, in contrast to its otherwise comprehensive STM design, transactions do
not support side-effects because they may be retried multiple times before
committing. Any side-effect performed within a transaction would therefore be
repeated on each retry.

\citeauthor{jensen2016}~\cite{jensen2016} addressed this by forking the
Clojure runtime and extending \texttt{LockingTransaction}, introducing
hooks such as \texttt{on-commit} and \texttt{after-commit}. While
effective, this approach couples the extension tightly to the runtime.

Rather than modifying \texttt{LockingTransaction}, we implement side-effect
support within our STM system at the language level. This avoids dependence
on a particular Clojure version while providing commit and rollback semantics.

Furthermore, Section~\ref{subsec:opacity-proof} shows that the resulting
system satisfies opacity~\cite{guerraoui2007}, providing formal correctness
guarantees for all transactions, including those that are aborted.